% Generated draft: requires mathematical and source review.
\documentclass[11pt]{article}
\usepackage[margin=25mm]{geometry}
\usepackage{fontspec}
\setmainfont{Latin Modern Roman}
\usepackage{amsmath,amssymb,mathtools,unicode-math}
\setmathfont{Latin Modern Math}
\usepackage{enumitem,needspace,xurl,hyperref}
\setlength{\parindent}{0pt}
\setlength{\parskip}{5pt}
\setlength{\emergencystretch}{3em}
\newcommand{\N}{\mathbb{N}}
\newcommand{\Z}{\mathbb{Z}}
\newcommand{\Q}{\mathbb{Q}}
\newcommand{\R}{\mathbb{R}}
\newcommand{\C}{\mathbb{C}}
\newcommand{\F}{\mathbb{F}}
\newcommand{\A}{\mathbb{A}}
\newcommand{\PP}{\mathbb{P}}
\newcommand{\E}{\mathbb{E}}
\newcommand{\eps}{\varepsilon}
\newcommand{\dd}{\,\mathrm{d}}
\begin{document}
\section*{UnsolvedMath ID: 11000130}
\textbf{Source problem code:} catalogue entry
\subsection{Definitions and mathematical statement}
For the usual unpointed genus-\(g\geq2\) convention, \(\mathcal M_g\) is the moduli stack of smooth compact complex curves and \(\overline{\mathcal M}_g\) its Deligne--Mumford compactification. Its added points represent connected stable nodal curves of arithmetic genus \(g\), with finite automorphism group. All stable dual graphs are allowed; this is not only the compact-type locus.

The question asks whether the source's Theorem 2.1 can be generalized from smooth-curve moduli to this compactification. The supplied summary describes the intended theorem as Madsen--Weiss-type, but the theorem's actual statement, homotopy/cohomology objects, maps, genus stabilization and range are absent. Accordingly the target is an analogue of that specific theorem for \(\overline{\mathcal M}_g\), with its precise objects and comparison supplied, rather than an invented infinite-loop-space equivalence. One must also specify stack/orbifold versus coarse-space interpretation and any marking conventions. The record gives no candidate formula or claim that a related model already proves the desired generalization.
\subsection{Short English statement}
Generalize the missing Theorem 2.1 to Deligne--Mumford stable-curve moduli, identifying the intended homotopy-theoretic comparison and boundary conventions.
\subsection{Sources}
\begin{itemize}
\item \emph{Problems on Mapping Class Groups and Related Topics}, edited by Farb, Question 4.4, supplied PDF page 106. \url{https://www.math.uchicago.edu/~farb/papers/mcgbook.pdf}. Recorded access 2026-07-29; extraction pdf-font-structure. Theorem 2.1 is not supplied.
\item The supplied summary identifies a Madsen--Weiss-type question and mentions related Ebert--Giansiracusa homotopy models of the stable-curve moduli stack, without complete bibliography entries or an exact resolving theorem.
\item Local UnsolvedMath registry, \texttt{data/unsolvedmath/problems.json}, permanent ID \(11000130\), exact source code \texttt{catalogue entry}; supplied OPEN-TRIAGE summary reports no complete resolution.
\end{itemize}
\end{document}
